% ---------------------------------------------------------------------------
\appendix

\section{Definitions of the main measures}
\label{app:measures}

Notation: a line is a sequence of words $w_1,\dots,w_n$; a word is a sequence of glyphs (Section~\ref{sec:data}); $\ell(w)$
and $f(w)$ are its last and first glyph. Plug-in entropies and mutual informations are in bits. Unless stated, a null is
averaged over 50--1{,}000 permutations and $z$ is the distance of the observed value from the null mean in null standard
deviations.

\paragraph{Junction} $I = I(\ell(w_i); f(w_{i+1}))$ over all pairs of adjacent words in the same line; $E = I -
\overline{I}_{\text{null}}$, the null shuffling the words of each line (same marginals and sample size, hence the same
finite-sample bias). The variant with position (\expcode{e3c87}) shuffles $f(w_{i+1})$ among pairs with the same class of
position (index of $w_i$ in the line: 0, 1, 2, 3, 4 or more; whether $w_{i+1}$ ends the line).

\paragraph{Mirroring of within-word rules} $\mathrm{PMI}_{\text{in}}(a,b) = \log \frac{n_{\text{in}}(a,b)\,N_{\text{in}}}
{n_{\text{in}}(a\cdot)\,n_{\text{in}}(\cdot b)}$ over consecutive glyphs inside words, and $\mathrm{PMI}_{\text{x}}(a,b)$
likewise over pairs $(\ell(w_i), f(w_{i+1}))$; $\rho$ is Spearman's correlation over the cells with at least five counts in
both tables.

\paragraph{Predictability of spaces} A rule $p(\text{space}\mid a,b)$ is estimated on the first half of the lines over
all pairs of consecutive glyphs $(a,b)$ (backing off to $a$, then to the global rate); on the second half a space is
predicted when $p > 0.5$; the score is the F1 of the predicted spaces.

\paragraph{Attested wrong cuts} The glyph stream of each line is re-spaced at random with the rule above (five
re-spacings); among the pieces of at least three glyphs that do not coincide with an original word, the share that are
word types of the text (median over re-spacings).

\paragraph{Form filling and frequency--form correlation} A first-order glyph chain with start and end symbols (add-0.5
smoothing) is estimated on the word types. For each length $L = 3,\dots,7$ the $K_L$ most probable sequences are found
by exact search, $K_L$ being the number of attested types of that length; filling is the share of them that are attested.
The frequency--form correlation is Spearman's correlation, within length classes and pooled after centring the ranks,
between the per-glyph log-probability of a type and the logarithm of its frequency (types occurring at least twice).

\paragraph{Agreement of choices (short state)} For a choice $c$ with forms coded $x \in \{0,1\}$ (for example \eva{k} = 1,
\eva{t} = 0), every inner word of a line of at least six words that shows the choice is an occurrence, with its
\emph{covered word} (the word with the variant position masked), hand, page and number of glyphs $g$ written before it in
the line. Its expected value is
\[
p_a = \bar x_{\text{hand,covered}}^{(-a)} + \big(\bar x_{\text{page}}^{(-a,-b)} - \bar x_{\text{hand}}^{(-a)}\big)
      + \beta_c\,\big(g_a - \bar g_{\text{page}}^{(-a,-b)}\big),
\]
clipped to $[0.01, 0.99]$, where the superscripts mean leave-one-out (or leave-two-out) means and $\beta_c$ is the
within-stratum slope of $x$ on $g$. For pairs $(a,b)$ of occurrences of the same choice at distance $d$ words in the same
line, excluding pairs whose covered words are equal or one edit apart,
\[
K(d) = \frac{\sum_{(a,b)} \big[\mathbb{1}(x_a = x_b) - e_{ab}\big]}{\sum_{(a,b)} \big[1 - e_{ab}\big]},\qquad
e_{ab} = p_a p_b + (1-p_a)(1-p_b).
\]
Because $p$ is estimated from the same data, $K$ is biased in a way that depends on the text (by up to $+0.37$). The
corrected value is $K_{\mathrm{c}}(d) = K(d) - \overline{K}(d)$, where $\overline{K}$ is the mean of $K$ over 50 shuffles of
the values $x$ among occurrences with the same hand, choice and covered word: the shuffles keep the bias and remove any
real agreement. On synthetic texts the correction brings a text without agreement back to zero and leaves a true
agreement unchanged (\expcode{e3c48}).

\paragraph{Drift along the line} Within strata (choice $\times$ covered word), the least-squares slope of $x$ on $g$,
pooled as $\sum \mathrm{cov} / \sum \mathrm{var}$ over strata, leaving out the first and last word of each line; reported
per ten glyphs.

\paragraph{Copying from the lines above} For paragraphs of at least four lines and lines from the third on, the share of
words that have an identical word, or one at one edit, in the two lines above; the null shuffles the order of the lines
within the paragraph, so that the ``lines above'' are other lines of the same paragraph; $E = $ observed $-$ null mean.

\paragraph{Left margin} For the lines of a paragraph from the second on, with at least three words and a first word of at
least two glyphs, the number of consecutive lines whose first two glyphs coincide, divided by its mean over 1{,}000
shuffles of the order of the lines within the paragraph.

\paragraph{Repetition} The share of adjacent pairs in a line that are the same word (words of at least two glyphs) or
differ by one edit (words of at least three glyphs), and its ratio to the mean share after shuffling the words within
each line (20 shuffles), or to $\sum_w p_w^2$.

\paragraph{Intervals} Contributions are summed per page; pages are grouped by bifolium (or quire), and groups are
resampled with replacement 2{,}000 times; intervals are the 2.5th and 97.5th percentiles. For the junction the null is
recomputed inside each resample from 20 within-line shuffles stored per page, so that observed and null share the same
finite-sample bias.

\begin{table}[htbp]
\centering\footnotesize
\caption{The main tests corrected for multiplicity: Holm within seven families, and Bonferroni over the 715 experiments
logged before this test ($p < 7.0\times10^{-5}$, $|z| > 3.97$). Experiment \expcode{e3c92}.}
\label{tab:multiplicity}
\begin{tabularx}{\linewidth}{>{\raggedright\arraybackslash}p{2.8cm}Yrcc}
\toprule
family & test & statistic & Holm & Bonferroni \\
\midrule
junction & junction between adjacent words & $z = 191.7$ & yes & yes \\
 & \eva{qo-}/\eva{o-} after the previous word & $z = 19.4$ & yes & yes \\
 & \eva{-l}/\eva{-r} before the next word & $z = 12.7$ & yes & yes \\
 & look-ahead, end of the first word changes & $p = 0.0002$ & yes & no \\
 & look-ahead, next word adapts & $p < 0.0001$ & yes & no \\
line edges & \eva{-m} at line end; \eva{-r} drops & $z = 26.7$; $-13.6$ & yes & yes \\
 & \eva{y-}, \eva{s-} at line start & $z = 16.6$; $16.0$ & yes & yes \\
copying & beyond the paragraph's vocabulary & $z = 17.4$ & yes & yes \\
left margin & \eva{qo-}, \eva{o-} avoided below the same start & $z = -9.7$; $-7.0$ & yes & yes \\
 & \eva{d-}, \eva{ch-}, \eva{y-} avoided & $z = -3.8$; $-3.6$; $-3.5$ & yes & no \\
short state & adjacent; two--three words (by bifolium) & $z = 12.3$; $8.1$ & yes & yes \\
drift & \eva{sh}/\eva{ch} & $z = -6.3$ & yes & yes \\
 & \eva{qo}/\eva{o}, \eva{k}/\eva{t}, \eva{-ey}/\eva{-dy} & $z = -3.0$; $-3.9$; $-2.8$ & yes & no \\
book & bifolium identity; languages A and B & $z = 6.6$; $13.4$ & yes & yes \\
\bottomrule
\end{tabularx}
\end{table}

% ---------------------------------------------------------------------------
\section{The log of experiments}
\label{app:log}

The repository contains one row per experiment from e01 to e3c92
(\url{white_paper/revisione/registro_esperimenti.csv}), with date, question, files, outcome against the
preregistered criterion, deviation and notes, and a summary with the borderline classifications. The log was compiled
semi-automatically from the preregistrations, results and notebook; every quoted outcome was checked to appear verbatim
in the sources, and 110 borderline classifications are listed for the reader to judge.

\begin{center}\footnotesize
\begin{tabular}{lrrrrrrrr}
\toprule
series & confirmed & refuted & mixed & inconclusive & superseded & descriptive & not run & total \\
\midrule
e01--e99 & 20 & 25 & 8 & 9 & 2 & 32 & 2 & 98 \\
e100--e199 & 10 & 61 & 9 & 18 & 9 & 3 & 2 & 112 \\
e200--e299 & 18 & 66 & 4 & 5 & 9 & 6 & 7 & 115 \\
e300--e399 & 37 & 26 & 15 & 13 & 2 & 7 & 1 & 101 \\
e3a & 43 & 24 & 12 & 2 & 12 & 4 & 2 & 99 \\
e3b & 22 & 22 & 12 & 12 & 29 & 1 & 1 & 99 \\
e3c & 20 & 35 & 11 & 13 & 6 & 4 & 3 & 92 \\
\midrule
all & 170 & 259 & 71 & 72 & 69 & 57 & 18 & 716 \\
\bottomrule
\end{tabular}
\end{center}

\noindent ``Confirmed'' and ``refuted'' refer to the affirmative answer to each preregistered question. Experiments
e01--e28 were imported from earlier work and re-run, without preregistration. ``Superseded'' marks outcomes later
withdrawn or reversed (in four cases only the method was corrected and the result upheld). Of the 80 declared deviations,
19 are outcomes read differently from the literal criterion; the most consequential is \expcode{e384}, whose criterion
(set from the range of short, noisy texts) literally gave ``the junction carries over the line break as in prose'',
while the measurements show a closed line ($Q \approx 0$, against 0.49 in languages; \expcode{e3c89}).

% ---------------------------------------------------------------------------
\section{Internal checks: provisional readings corrected by later work}
\label{app:checks}

\paragraph{During the research (before version 1)}
\begin{enumerate}
\item The agreement measure was biased by up to $+0.37$ (Appendix~\ref{app:measures}). Before the correction it had
  given the Anglo-Saxon scribe of the \emph{Hatton Gospels} a ``window'' like the Voynich's, made Volapük look similar to
  the Voynich, placed the shape of the Voynich's agreement outside the range of languages (it is at the top), and suggested
  that the state travels only on occurrences of the choice.
\item The state of the choices first looked like a window of three words dropping to zero; it fades gradually (half-life
  about three words), and a slow component crosses lines.
\item A per-line ``setting'' of the choices, like a key per line, did not hold.
\item After the form-only variants of the Glen Claston transliteration we read the state as tied to words rather than to
  letter shapes; a real scribe with a window in a form-only choice made that generalisation too broad.
\item An uncorrected measure suggested that hand 1 lacks the state; corrected, all hands have it.
\item An apparent reset of the choices at the line break came from comparing with the page average (line edges have
  forms of their own).
\item The first left-margin test used a defective null; the corrected one confirms the result.
\item A share of the link between whole words attributed to the paragraph was a mechanical effect of smaller groups.
\item Rare forms seemed to be copied from the line above more than others; with the second transliteration only a trend
  remains.
\item Two formally positive reading candidates (a syllabic one and a German one) were artefacts: the first came out the
  same on the shuffled Voynich, the second changed sign with the starting points of the search.
\item Three criteria called ``absent'' an effect whose interval only touched zero; the estimates were positive and
  similar to the rest. Since then ``absent'' requires the upper end of the interval below half the reference value.
\end{enumerate}

\paragraph{Prompted by the external review of version 1}
\begin{enumerate}[resume]
\item ``71 languages'' were 71 texts, 15 of them in constructed languages; per natural language almost all comparisons
  hold, but two maxima came from Toki Pona, the junction is exceeded by two languages of 24 (not ``six of 71''), and the
  per-text junction values of version~1 came from an unsaved exploratory check with a different glyph segmentation
  (\expcode{e3c84}).
\item The ``weak space inside a chain of glyphs'' does not survive a word-grammar null: predictable spaces, attested wrong
  cuts and form filling follow from word form (\expcode{e3c85}). The junction and its mirroring of within-word rules
  survive.
\item ``Repeats the previous word more than chance'' becomes ``does not avoid repeating it, and repeated words cluster in
  the line'' (\expcode{e3c88}); ``no comparison text has this signature'' for the decline of repetition within the line was
  wrong (two languages of 20; \expcode{e3c84}).
\item The drift was said to be ``three to fifteen times'' that of the Nordic scribes; relative to the strongest scribal
  drift it is about 2.5 to 4.5 times.
\item ``Not found in any of 79 scribes'' was too strong for the short state: with adequate power it is absent in 14 of
  29 choices, and Holm A 10 shows a weaker state-like agreement (\expcode{e3c90}, \expcode{e3c91}). In the first run of
  \expcode{e3c90} the drift intervals of the Nordic scribes were read in the wrong unit (per glyph instead of per ten
  glyphs), giving 28 instead of 20 choices ``absent with power''; corrected the same day.
\item Version~1 described an early version of the voynichizer (message in the spelling choices); the published version
  carries it in the word counts.
\item Version~1 presented as new several regularities described before (Section~\ref{sec:related}).
\end{enumerate}

% ---------------------------------------------------------------------------
\section{Glossary}

\begin{description}[leftmargin=0pt,itemsep=2pt]
\item[EVA] conventional alphabet to transliterate the Voynich into Latin letters; it does not indicate sounds.
\item[glyph] a sign of the manuscript; \eva{ch}, \eva{sh} and the bench gallows count as one.
\item[gallows] the tall signs \eva{k}, \eva{t}, \eva{p}, \eva{f}.
\item[Currier languages A and B] the two varieties of writing into which the pages divide.
\item[junction] information that the last glyph of a word gives about the first glyph of the next, beyond chance.
\item[junction rule] a rule linking the form of a word to its neighbour (sandhi-like).
\item[variant forms] pairs of words that differ by one glyph in a fixed position (for example \eva{k}/\eva{t}).
\item[short state] the tendency to repeat the same variant for a few words.
\item[scorecard] the 18 properties with pass bands that the generator must reproduce.
\item[judge] a classifier retrained on each generated book to tell its pages from the real ones; its AUC is a classifier
  two-sample test.
\item[AUC] area under the ROC curve: 0.5 for guessing, 1 for perfect separation.
\item[arithmetic decoder] the inverse of an arithmetic coder; fed with random bits it samples from a probability model.
\item[preregistration] a dated description, written before running, of question, data, measure and criterion.
\item[95\% interval] a bootstrap percentile interval: the range of the middle 95\% of the estimates over resamples of
  bifolia (or pages).
\end{description}

% ---------------------------------------------------------------------------
\section{Data sources and licences}
\label{app:sources}

\begin{center}\footnotesize
\begin{tabularx}{\linewidth}{>{\raggedright\arraybackslash}p{4.5cm}Yp{3.2cm}}
\toprule
data & source & licence \\
\midrule
Voynich transliterations ZL3b, IT2a, GC2a & \citet{zandbergen_transliterations} & public domain (CC0); in the repository \\
71 meaningful texts, 38 gibberish texts & \citet{gaskellbowern2022}, released subset & modified MIT (cite the paper) \\
Bible in 100 languages & \citet{christodouloupoulos2015} & CC0 \\
Latin Library; Divinum Officium & CLTK collection; Divinum Officium project & as stated by the sources (MIT for Divinum Officium) \\
six Nordic manuscripts & \citet{menota}; each text with the editor and version in its header & CC BY-SA 4.0 \\
73 German manuscripts & \citet{ref2021}, v1.0.2 & CC BY 4.0 (Zenodo) or CC BY-SA 4.0 (file); we follow the latter \\
Copiale & \citet{knight2011,knight2012}, project page at Stockholm University & research use \\
Naibbe & \citet{greshko2025}, repository greshko/naibbe-cipher & modified MIT \\
Timm and Schinner generator & \citet{timmschinner2020} & MIT \\
U2, U3 & \citet{whitehatnetizen2026} & MIT \\
readings tested & \citet{bax2014,vatne2021,cheshire2019,schechter2026,gatta2026} & as published; Schechter: none stated, local copy not redistributed \\
\bottomrule
\end{tabularx}
\end{center}
